\par\Needspace{6\baselineskip}\medskip
\hypertarget{native-083-article064}{}\pdfbookmark[2]{64.}{article-064}
\begin{center}64.\end{center}
Zu der Aufgabe aus Artikel 62 fügen wir noch die folgenden Bemerkungen hinzu.

I. Setzt man in der dort angegebenen zweiten Gleichung $N=\lambda$, $N=L$, $N=l$, so erhält man $R'\sin(\lambda-L)=\Delta'\sin(l-\lambda)$, $r'\sin(\lambda-L)=\Delta'\sin(l-L)$, $r'\sin(l-\lambda)=R'\sin(l-L)$. Die erste oder die zweite Gleichung eignet sich zur Kontrolle der Rechnung, wenn man Methode I oder II aus Artikel 62 verwendet hat. In unserem Beispiel erhält man damit
\begin{center}
\begin{tabular}{@{}ll@{}}
\makebox[13em][l]{$\log\sin(\lambda-L)$\hfill .\,.\,.} & $9{,}475\,8653_n$\\[.3em]
\makebox[13em][l]{$\log\frac{\Delta'}{r'}$\hfill .\,.\,.\,.\,.\,.\,.\,.\,.} & $9{,}754\,6117$\\[.3em]
\hline
& $9{,}721\,2536_n$\\[.3em]
\makebox[13em][l]{$\log\sin(l-L)$\hfill .\,.\,.} & $9{,}721\,2536_n$
\end{tabular}
\end{center}
\begin{center}
$l-L=-31^0\,45^{\prime}\,26^{\prime\prime}{,}82$
\end{center}

II. Die Sonne und die beiden Punkte in der Ebene der Ekliptik, auf die der Ort des Himmelskörpers und der Ort der Erde projiziert werden, bilden ein ebenes Dreieck mit den Seiten $\Delta'$, $R'$, $r'$ und den gegenüberliegenden Winkeln entweder $\lambda-L$, $l-\lambda$, $180^0-l+L$ oder $L-\lambda$, $\lambda-l$, $180^0-L+l$. Daraus folgen unmittelbar die in I angegebenen Beziehungen.

III. Die Sonne, der wahre Ort des Himmelskörpers im Raum und der wahre Ort der Erde bilden ein weiteres Dreieck mit den Seiten $\Delta$, $R$, $r$. Bezeichnen wir die jeweils gegenüberliegenden Winkel mit $S$, $T$, $180^0-S-T$, so gilt $\frac{\sin S}{\Delta}=\frac{\sin T}{R}=\frac{\sin(S+T)}{r}\cdot$ Die Ebene dieses Dreiecks ergibt auf der Himmelskugel einen Großkreis. Auf ihm liegen der heliozentrische Ort der Erde, der heliozentrische Ort des Himmelskörpers und dessen geozentrischer Ort so, dass der Abstand des zweiten vom ersten, der des dritten vom zweiten und der des dritten vom ersten, in derselben Richtung gezählt, jeweils $S$, $T$, $S+T$ betragen.

\SourceLeaf{084}IV. Die folgenden Differenzialgleichungen lassen sich aus den bekannten differenziellen Änderungen der Stücke eines ebenen Dreiecks oder ebenso leicht aus den Formeln von Artikel 62 herleiten.
\[
\mathrm{d}l=\frac{r'\cos(\lambda-l)}{\Delta'}\,\mathrm{d}\lambda+\frac{\sin(\lambda-l)}{\Delta'}\,\mathrm{d}r'
\]
\[
\mathrm{d}\Delta'=-r'\sin(\lambda-l)\,\mathrm{d}\lambda+\cos(\lambda-l)\,\mathrm{d}r'
\]
\[
\begin{aligned}
\mathrm{d}b={}&\frac{r'\cos b\sin b\sin(\lambda-l)}{\Delta'}\,\mathrm{d}\lambda
+\frac{r'\cos b^2}{\Delta'\cos\beta^2}\,\mathrm{d}\beta\\
&+\frac{\cos b^2}{\Delta'}(\operatorname{tang}\beta-\cos(\lambda-l)\operatorname{tang}b)\,\mathrm{d}r',
\end{aligned}
\]
\noindent Werden die Winkeländerungen in Bogensekunden ausgedrückt, so sind hier die Terme, die $\mathrm{d}r'$, $\mathrm{d}\Delta'$ enthalten, mit $206265^{\prime\prime}$ zu multiplizieren oder die übrigen durch $206265^{\prime\prime}$ zu teilen.

V. Die umgekehrte Aufgabe, den heliozentrischen Ort aus dem geozentrischen zu bestimmen, ist der oben behandelten ganz analog. Es wäre daher überflüssig, länger bei ihr zu verweilen. Sämtliche Formeln aus Artikel 62 gelten auch für diese Aufgabe, wenn man alle auf den heliozentrischen Ort des Himmelskörpers bezogenen Größen mit den entsprechenden Größen für den geozentrischen Ort vertauscht und für $L$, $B$ entsprechend $L+180^0$, $-B$ einsetzt. Gleichbedeutend damit ist es, den geozentrischen Ort der Sonne anstelle des heliozentrischen Ortes der Erde zu nehmen.
